\subsection{Ethics and Broader Impact}

This work focuses on stablecoin arbitrage across centralized exchanges (CEXs), operating within the established legal and regulatory frameworks of major cryptocurrency platforms. All data used in our experiments are publicly available market prices and order-book snapshots accessible through standard exchange APIs. No proprietary, private, or user-specific data are used.

\paragraph{Market impact.}
Arbitrage plays a well-documented role in financial markets by improving price efficiency and reducing cross-venue price discrepancies. In stablecoin markets specifically, arbitrage trading helps maintain the dollar peg by correcting mispricings across exchanges. Our system operates at portfolio sizes (\$10,000--\$100,000) that are unlikely to move market prices on major exchanges, and the 3-hop path structure limits the number of concurrent market interactions. We note, however, that widespread deployment of high-frequency arbitrage systems could increase competition for these opportunities, potentially compressing margins and increasing latency sensitivity.

\paragraph{Execution risk and financial loss.}
Our framework is a research prototype and has not been deployed for live trading with real capital. The results presented represent \emph{planned} arbitrage paths based on real-time market snapshots, not executed trades. As discussed in Section~6, order books can change between planning and execution, and the profitability of a planned path is not guaranteed at execution time. Users adapting this system for live trading should incorporate additional safeguards including position limits, circuit breakers, and real-time slippage monitoring.

\paragraph{Environmental considerations.}
Our system relies on centralized exchange APIs and does not directly interact with blockchain networks during the search phase. Consequently, it does not contribute to blockchain energy consumption beyond standard API usage. However, executing the discovered arbitrage paths would require on-chain transfers, contributing to network congestion and energy use proportional to the chosen blockchain's consensus mechanism.

\paragraph{Regulatory compliance.}
Cryptocurrency arbitrage is legal in most jurisdictions, though regulatory environments vary. Users should ensure compliance with local financial regulations, including know-your-customer (KYC) requirements, tax reporting obligations, and any applicable trading restrictions. Our research does not constitute financial advice.

